\chapterimage{golden_gate.jpg}

\chapter{Capa 3: comunicació entre xarxes (Internet)}\label{sec:layer3}

\begin{minipage}{0.4\linewidth}
\begin{center}
\begin{bytefield}{16}
\bitbox{16}{Layer 7: Application} \\
\bitbox{16}{Layer 4: Transport} \\
\bitbox{16}{\color{color1} Layer 3: Internet} \\
\bitbox{16}{Layer 2: Network (LAN)} \\
\bitbox{16}{Layer 1: Physical} \\
\end{bytefield}
\end{center}
\end{minipage}
\begin{minipage}{0.6\linewidth}
\begin{center}
\includegraphics[width=\linewidth]{network_layer3.pdf}
\end{center}
\end{minipage}

\vspace{-0.75cm}


% \section{Layer~3 overview}
\subsection*{Funcionalitats}

El \conceptRef{IP}{protocol d'Internet (IP)} permet intercanviar \conceptRef{datagram}{datagrames} 
(\conceptRef{PDU}{PDU} de la capa~3) entre dispositius de diferents \conceptRef{LAN}{LANs} (xarxa d'àrea local).
% 
El seu sistema d'adreçament identifica els dispositius globalment i fa possible 
\conceptRef{routing}{encaminar} els datagrames allà on calgui. 

Els dispositius que es comuniquen via IP poden pertànyer a LANs
basades en tecnologies diferents i amb tipus d'adreça \concept{MAC} diferents.
Això té dues implicacions:
\begin{itemize}
  \item Un datagrama que ``cap'' en la \concept{MTU} de la LAN d'origen
    pot ser massa gran per a la MTU de la LAN de destinació (o de qualsevol LAN intermèdia). 
    IP defineix un sistema de \conceptRef{fragmentation}{fragmentació} per dividir els datagrames,
    que només es reconstrueixen a la destinació final.\\[-0.25cm]
    
  \item Quan enviem un datagrama, coneixem l'\conceptRef{IP}{adreça IP} del dispositiu destinatari,
  però no necessàriament la seva adreça \concept{MAC}. Només necessitarem aquesta MAC si formem part de la LAN
  de destinació. IP fa servir \concept{ARP} per traduir adreces IP conegudes
  a adreces MAC del tipus de la LAN de destinació.
\end{itemize}


Els \conceptRef{datagram}{datagrames} poden viatjar per xarxes que no controlen 
ni l'origen ni la destinació, de manera que els datagrames es poden perdre, desordenar i duplicar.
% 
La capa~3 no proporciona mecanismes per gestionar-ho; les capes superiors han d'implementar
mecanismes de protecció quan calgui 
(\eg, per implementar els \conceptRef{stream}{fluxos de dades} de \concept{TCP}).

\begin{center}
\includegraphics[width=0.9\linewidth]{protocols_layer3.pdf}
\end{center}


\vspace{-0.25cm}
\subsection*{Protocols}

La versió~4 d'\concept{IP} és \textit{el} protocol que s'utilitza a la capa~3, \ie,
és comú a pràcticament totes les comunicacions de la Internet tal com la coneixem avui.
% 
Fa molts anys que Internet lluita per fer la transició cap a \concept{IPv6}, 
però el procés no s'ha completat i només \concept{IPv4} ofereix cobertura global.

\concept{IPv4} delega en el protocol \concept{ARP} la traducció d'adreces IP conegudes 
a adreces MAC dins de cada LAN. ARP \textit{no} fa servir datagrames IP.

Els protocols següents encapsulen les seves \conceptRef{PDU}{PDU} en el payload (càrrega útil) dels 
\conceptRef{datagram}{datagrames} IP:\\[-0.6cm]
\begin{itemize}
\item Transport Control Protocol (\concept{TCP}) 
\item User Datagram Protocol (\concept{UDP})
\item Internet Control Message Protocol (\concept{ICMP})
\end{itemize}

\begin{remark}
Les ordres \inlineCode{ip address}, \inlineCode{ip route} i \inlineCode{ip neighbor} 
(entre d'altres) et permeten consultar diversos aspectes de la teva configuració IP.
\end{remark}



\section{Adreçament de la capa~3}\label{sec:layer3:addressing}

Les adreces \concept{IPv4} tenen $32$~bits. Normalment es presenten a les persones en
el format \concept{quad decimal}, \eg,
\begin{center}
\otherBase{142.250.201.67}
\end{center}

Les adreces \concept{IPv6} tenen $128$~bits ($12$~bytes), \ie, 
són $4$~vegades més grans que una adreça IPv4.
S'expressen en format hexadecimal, \eg, com es mostra a continuació. Fixa't que 
dos signes de dos punts ``\otherBase{::}'' indiquen ``omple amb zeros'',
i que ``\otherBase{::}'' només pot aparèixer una vegada per adreça IPv6.
\begin{center}
\otherBase{fe80::f524:a73a:e946:7c02}
\end{center}

\begin{remark}
Quan una adreça \concept{IPv6} forma part d'un URL, s'envolta de claudàtors,
\eg, \otherBase{https://[fe80::f524:a73a:e946:7c02]/}.
\end{remark}


\subsection{Adreces públiques i reservades}

La majoria d'adreces IP són \conceptRef{public IP}{públiques} i úniques a tot Internet, 
\ie, l'adreça IP anterior té el mateix significat a tot el món: identifica un únic dispositiu connectat.
% 
Moltes adreces són \conceptRef{reserved IP address}{reservades/privades} i només es poden 
fer servir dins d'un àmbit concret (\eg, una LAN) o en una situació concreta.
Algunes de les més habituals són les següents (n'hi ha 
\href{https://en.wikipedia.org/wiki/Reserved_IP_addresses#IPv4}{\underline{algunes més}}, 
també per a IPv6):

\begin{center}
\begin{tabular}{ll}
\toprule
\textbf{Bloc d'adreces IP} & \textbf{Àmbit de l'adreça} \\
\toprule
\otherBase{127.*.*.*} & Aquest ordinador (localhost) \\[0.1cm]
\otherBase{10.*.*.*} & Aquesta LAN \\
\otherBase{172.16.*.*} & Aquesta LAN \\
\otherBase{192.168.*.*} & Aquesta LAN \\[0.1cm]
\otherBase{0.*.*.*} & Usos especials \\
\otherBase{169.254.*.*}& Usos especials \\
\otherBase{255.255.255.255} & Usos especials\\
\bottomrule
\end{tabular}
\end{center}

\begin{exercise}\ \\[-0.5cm]
\begin{itemize}
\item Quantes adreces \concept{IPv4} diferents hi ha (incloent-hi les reservades i privades)? 
\item Són suficients ara i en el futur?
\item És problemàtic que una adreça com \otherBase{192.168.0.1} 
  l'estiguin fent servir simultàniament milers (probablement milions) d'ordinadors a Internet ara mateix?
\item Quina d'aquestes correspon a adreces \concept{link local}?
\end{itemize}
\end{exercise}

\subsection{Màscares de xarxa}\label{sec:layer3:netmasks}

Els $32$~bits d'una adreça IP es poden dividir en dues parts mitjançant una \conceptRef{netmask}{màscara de xarxa}:\\[-0.25cm]
% 
\begin{center}
\begin{bytefield}{32}
\bitheader{0-31}\\
\bitbox{10}{Network ID} & \bitbox{22}{Device ID} \\
\end{bytefield}
\end{center}
% 
La primera part identifica la LAN o el grup de LANs a què pertany un dispositiu. 
La segona part identifica un únic dispositiu dins de la LAN. 
La màscara de xarxa determina quants bits s'assignen a cada part. 

A la figura anterior, la màscara de xarxa assigna 10 bits a l'identificador de xarxa i 22 al del dispositiu.
Aquesta màscara pot anar des de $0$ (tots els bits són identificador de dispositiu) fins a $32$ (tots els bits són identificador de xarxa).


Les màscares de xarxa es presenten a les persones de dues maneres principals:\\[-0.5cm]
\begin{itemize}
  \item \otherBase{/N}, on $N$ és el nombre de bits assignats a l'identificador de xarxa. 
    És el que es fa servir en el format \concept{CIDR}.
  
  \item El format \concept{quad decimal} (\otherBase{A.B.C.D}) de 
    $N$ bits a \otherBase{1}, seguits de $(32-N)$ bits a \otherBase{0}.
\end{itemize}

\begin{exercise}
La màscara de xarxa de l'exemple anterior es pot expressar en binari 
(tot i que no és gaire pràctic):
\begin{center}
\begin{bytefield}[bitwidth=1em]{32}
\bitheader{0-31} \\
\bitboxes{1}{1111111111} &
\bitboxes{1}{0000000000 0000000000 00} \\
\bitbox{10}{$10$~bits} & \bitbox{22}{$32-N = 22$~bits}
\end{bytefield}
\end{center}

\begin{itemize}
\item Expressa la màscara de xarxa de l'exemple anterior en els formats CIDR i quad decimal.
\item Quina és la teva IP i la teva màscara de xarxa? Expressa-les de les mateixes dues maneres.
\end{itemize}
\end{exercise}

\begin{remark}
Les màscares de xarxa també es fan servir a IPv6, però amb un valor màxim de $N$ de $128$ en lloc de $32$.
\end{remark}

\subsection{Blocs d'adreces IP}

Les màscares de xarxa es combinen amb adreces IP per definir 
\conceptRef{IP address block}{blocs d'adreces IP} 
que comencen pels mateixos $N$ bits (on $N$ és la màscara de xarxa).
% 
Per exemple, \otherBase{192.168.3.0/24} (en format \concept{CIDR}) 
i \otherBase{192.168.3.0/255.255.255.0} fan referència al bloc 
d'adreces entre \otherBase{192.168.3.0} i \otherBase{192.168.3.255},
ambdues incloses
(\otherBase{192.168.3.*} en la notació informal de la taula anterior).

\begin{exercise} 
Els blocs IP poden ser més petits, iguals o més grans que una xarxa IP.
Dona exemples en què un bloc IP faci referència a:
\begin{itemize}
\item Una única IP.
\item La meitat de les adreces IP de la xarxa \otherBase{192.168.3.0/24}.
\item Totes les adreces de la xarxa \otherBase{192.168.3.0/24}, més 
      $256$ adreces més.
\item Totes les IP possibles.
\end{itemize}
\end{exercise}


\begin{exercise}\ \\[-0.65cm]
\begin{itemize}
\item Es pot expressar qualsevol \conceptRef{netmask}{màscara de xarxa} com una \conceptRef{IP address}{adreça IP}?
\item Enumera totes les màscares de xarxa possibles.\\[-0.5cm]
\end{itemize}
\end{exercise}

Quan es fa referència a un bloc d'adreces (com una LAN), 
l'adreça IP d'aquesta combinació IP + màscara de xarxa
ha de tenir tots els bits de l'identificador de dispositiu a \otherBase{0}. 
D'aquesta manera, \otherBase{192.168.3.0/24} fa referència correctament a una xarxa, però
\otherBase{192.168.3.1/24} fa referència, en canvi, a un dispositiu 
(el que té la IP \otherBase{192.168.3.1}) d'aquesta xarxa (\otherBase{192.168.3.0/24}).

\subsection{LANs i màscares de xarxa}\label{sec:layer3:lans_and_netmasks}

Les màscares de xarxa \textit{no} han de ser necessàriament múltiples de $8$. Considera l'exemple següent, en què \\
es fa servir \otherBase{192.168.76.1/19} (la IP + màscara d'un dispositiu) per obtenir la seva adreça de xarxa\\
(\otherBase{192.168.64.0/19}):\\

\begin{center}
\begin{bytefield}[bitwidth=1.1em]{32}
\bitheader{0-31}\\
\begin{rightwordgroup}{Device's IP}
\bitbox{8}{192} & \bitbox{8}{168} & \bitbox{8}{76} &\bitbox{8}{1} \\
\bitboxes{1}{11000000} & \bitboxes{1}{10101000} & \bitboxes{1}{0100{\bf1}{\bf1}00} & \bitboxes{1}{0000000{\bf1}}
\end{rightwordgroup} \\

\bitboxes[]{1}{\&\&\&\& \&\&\&\& \&\&\&\& \&\&\&\& \&\&\&
               {\color{color1}\&} 
               {\color{color1}\&}{\color{color1}\&}{\color{color1}\&}{\color{color1}\&}
               {\color{color1}\&}{\color{color1}\&}{\color{color1}\&}{\color{color1}\&}
               {\color{color1}\&}{\color{color1}\&}{\color{color1}\&}{\color{color1}\&}} \\

\begin{rightwordgroup}{Netmask}
\bitboxes{1}{1111111111111111111
   {\color{color1}0}{\color{color1}0}{\color{color1}0}{\color{color1}0}
   {\color{color1}0}{\color{color1}0}{\color{color1}0}{\color{color1}0}
   {\color{color1}0}{\color{color1}0}{\color{color1}0}{\color{color1}0}
   {\color{color1}0}}
\end{rightwordgroup} \\

\bitboxes[]{1}{==== ==== ==== ==== ===
               {\color{color1}=} 
               {\color{color1}=}{\color{color1}=}{\color{color1}=}{\color{color1}=}
               {\color{color1}=}{\color{color1}=}{\color{color1}=}{\color{color1}=}
               {\color{color1}=}{\color{color1}=}{\color{color1}=}{\color{color1}=}} \\

\begin{rightwordgroup}{Network IP}
\bitboxes{1}{11000000} & \bitboxes{1}{10101000} & \bitboxes{1}{0100{\bf0}{\bf0}00} & \bitboxes{1}{0000000{\bf0}} \\
\bitbox{8}{192} & \bitbox{8}{168} & \bitbox{8}{64} &\bitbox{8}{0} 
\end{rightwordgroup}\\
\end{bytefield}
\end{center}

\begin{exercise}\ \\[-0.5cm]
\begin{itemize}
\item A l'exemple anterior, en realitat només calia transformar un valor decimal a binari,
i un valor decimal a binari. Quin?
\item Això és sempre cert quan es converteixen combinacions IP+màscara de xarxa a adreces de xarxa?
\item Expressa la màscara de xarxa en format \concept{quad decimal}.
\end{itemize}
\end{exercise}



Donat un bloc d'adreces IP + màscara de xarxa i l'adreça IP d'un dispositiu, és directe determinar 
si la IP d'aquest dispositiu pertany al bloc (si i només si coincideixen els primers $N$ bits,
on \otherBase{/N} és la màscara de xarxa CIDR).

\begin{exercise}
Determina si cadascuna de les adreces següents pertany a la mateixa xarxa que
el dispositiu \otherBase{123.45.67.89/14}:
\begin{itemize}
\item \otherBase{23.45.67.89}
\item \otherBase{123.45.203.25}
\item \otherBase{123.43.255.255}
\item \otherBase{0.45.67.1}
\end{itemize}
\end{exercise}

Una \conceptRef{network}{xarxa IP} (una LAN) és un bloc d'adreces, però no totes aquestes adreces 
es poden assignar a dispositius. Sempre hi ha dues adreces reservades:
\begin{itemize}
\item L'adreça amb tots els bits de node a \otherBase{0}: identifica la mateixa LAN.
\item L'adreça amb tots els bits de node a \otherBase{1}: identifica el \concept{broadcast} (difusió) dins de la LAN.
\end{itemize}

\begin{exercise}\ \\[-0.5cm]
\begin{itemize}
\item Enumera totes les adreces que es poden assignar a dispositius a \otherBase{192.168.54.0/29}.
\item Indica l'adreça de broadcast d'aquesta xarxa.
\end{itemize}
\end{exercise}

La màscara de xarxa \otherBase{/N} determina el nombre de dispositius que es poden connectar alhora
a una xarxa: $2^{32-N} - 2$. Quan es diu que una LAN és ``petita'' o ``més gran'', la mida fa referència a aquest 
nombre d'adreces de dispositiu disponibles. 

A la inversa, el nombre $M$ de dispositius que necessitem connectar a la mateixa LAN determina la 
mida de LAN més petita (la màscara \otherBase{/N} més gran) on caben tots aquests dispositius alhora:

\begin{center}
\begin{tabular}{ccc}
\toprule
\textbf{Màscara} & \textbf{Total adreces} & \textbf{Màx. dispositius} \\
\otherBase{/N} & $2^{32-N}$ & $2^{32-N} - 2$ \\
\toprule
\otherBase{/30} & 4 & 2 \\
\otherBase{/29} & 8 & 6 \\
\otherBase{/28} & 16 & 14 \\
\otherBase{/27} & 32 & 30 \\
\otherBase{/26} & 64 & 62 \\
\otherBase{/25} & 128 & 126 \\
$\vdots$ & $\vdots$ & $\vdots$ \\
\bottomrule
\end{tabular}
\end{center}

\begin{remark}
Tots els \conceptRef{router}{routers} (encaminador) necessiten una adreça IP vàlida per cada LAN a la qual estan connectats.
Per tant, els routers \textit{sí} que compten per al nombre màxim de dispositius connectats a una LAN.
\end{remark}

\begin{exercise} \ \\[-0.5cm]
\begin{itemize}
\item Quina és la mida de xarxa mínima on caben $256$ dispositius connectats?
\item Quina és la màscara \otherBase{/N} més llarga que permetria connectar-se a tothom de la teva ciutat?
\item Quantes adreces conté una xarxa \otherBase{/0}? 
\item Quantes xarxes \otherBase{/0} poden existir? Tenen nom?
\end{itemize}
\end{exercise}


\begin{exercise}
El fragment de codi següent pren una adreça IP i en mostra els $4$~bytes expressats en 
quad decimal i en binari.
Amplia (o substitueix) el codi perquè puguis:
\begin{itemize}
\item Transformar $32$ bits en una cadena quad decimal.
\item Mostrar una màscara de xarxa en format d'adreça IP i en format CIDR.
\item Donades la IP i la màscara de xarxa d'un dispositiu, determinar si una altra IP pertany a la mateixa xarxa.
\end{itemize}
\begin{center}
\showCode{snippets/ipcalculator.py}
\end{center}
\end{exercise}




% \todo{exercise netmask calculator}
% 
% \todo{HERE HERE HERE: add and point to the subnetting area}







\section{IP -- Internet Protocol}\label{sec:layer3:ip}

\subsection{Format del paquet}\label{sec:layer3:ip:format}

\concept{IP} defineix un format de \conceptRef{packet}{paquet} amb una \conceptRef{header}{capçalera} 
que té \textit{com a mínim} $20$~bytes. S'hi poden incloure parts opcionals,
sempre que la mida total de la capçalera sigui múltiple de $32$~bits ($4$~bytes).
El payload pot ser buit, de manera que la mida mínima d'un datagrama és de $20$~bytes.\\

\begin{center}
\small
\begin{bytefield}[bitheight=3.4em,bitwidth=1.1em]{32}
\bitheader{0-31}\\
\begin{rightwordgroup}[curlyshrinkage=10pt]{Mandatory Header\\($20$~bytes)}
  \bitbox{4}{IP\\Version} 
  & \bitbox{4}{IHL\\{\scriptsize(HLen/$4$)}} 
  & \bitbox{8}{Type of Service}
  & \bitbox{16}{Datagram Length\\{\scriptsize(Head + Payload, bytes)}} \\
  % 
  \bitbox{16}{Fragment ID} & \bitbox{1}{\otherBase{0}}
  & \bitbox{1}{\rotatebox{90}{\scriptsize DF flag}}
  & \bitbox{1}{\rotatebox{90}{\scriptsize MF flag}}
  & \bitbox{13}{Fragment Offset bytes/$8$} \\
  % 
  \bitbox{8}{Time to Live\\(TTL)}
  & \bitbox{8}{Payload protocol}
  & \bitbox{16}{IP Header Checksum}\\
  % 
  \bitbox{32}{Source IP Address} \\
  \bitbox{32}{Destination IP Address}
\end{rightwordgroup} \\
%   
\begin{rightwordgroup}[curlyshrinkage=10pt]{Optional Header\\{\scriptsize($4$ byte words)}}
\bitbox[tlr]{32}{Options\\{\scriptsize(variable length, optional)}} \\
\bitbox[lbr]{24}{} &
\bitbox{8}{Padding to\\$32$ bits}
\end{rightwordgroup} \\
% 
\begin{rightwordgroup}[curlyshrinkage=10pt]{Payload}
\bitbox[tlr]{32}[bitheight=3em]{Payload data\\{\scriptsize(variable length, byte aligned)}} \\
\bitbox[lbr]{24}{\hspace{3cm}\raisebox{0.25cm}{$\vdots$}} & \bitbox[lt]{8}{}
\end{rightwordgroup}
\end{bytefield}
\end{center}

\begin{exercise} \ \\[-0.5cm]
\begin{itemize}
\item Com es pot calcular la longitud del payload a partir de la capçalera?
\item Quina és la longitud màxima del payload?
\item Quins són els valors possibles dels flags (indicador) de la capçalera IP?
\item Què és més llarg, una adreça IP o una adreça MAC Ethernet?
\item És possible enviar exactament 17 bits de dades de payload en un datagrama IP?
\item Com podem saber si un datagrama encapsula dades TCP o UDP?
\item Per què la part opcional de la capçalera ha de ser múltiple de $32$~bits?
\end{itemize}
\end{exercise}

\subsection{Funcionament}

Quan un dispositiu A vol enviar un datagrama IP a un dispositiu B, les IP d'origen i 
de destinació de la capçalera queden fixades i no canvien per a aquest datagrama.

Si els dispositius A i B pertanyen a la mateixa LAN, el datagrama s'envia directament
(\eg, encapsulat en un frame (trama) Ethernet). Tanmateix, IP pot arribar molt més lluny.
Considera l'escenari següent, en què A vol comunicar-se amb E.

\begin{center}
 \includegraphics[width=0.5\linewidth]{ip_operation_example.pdf}
\end{center}


Aquí, els dispositius d'origen i de destinació pertanyen a LANs diferents.
Primer, el datagrama es passa al \concept{router} de l'origen.
Després, cada router intermedi el passa al següent fins que el datagrama arriba 
a l'últim router, que està connectat a la LAN de destinació. Finalment, aquest últim router 
passa el datagrama a la destinació. 

\begin{center}
\includegraphics{ip_operation.eps}
\end{center}

Les adreces IP d'origen i de destinació no canvien al llarg del trajecte.
En canvi, l'adreça MAC canvia a cada \conceptRef{hop}{salt} entre LANs. 

\section{Encaminament IP}\label{sec:layer3:routing}

Els \conceptRef{router}{routers} són els dispositius encarregats de reenviar els \conceptRef{datagram}{datagrames} IP
mirant l'adreça IP de destinació de cadascun, present a la \conceptRef{header}{capçalera} IP. Els routers no necessiten 
\conceptRef{decapsulate}{desencapsular} les capçaleres de les capes 4 i superiors. Per tant, les seves piles de xarxa només necessiten arribar fins a la capa~3
(i per això de vegades se'ls anomena dispositius de capa~3).

\begin{center}
\includegraphics[width=0.9\linewidth]{forward_layers.pdf}
\end{center}


Quan un router detecta un frame de xarxa $F$ en una de les seves \conceptRef{interface}{interfícies} (targetes de xarxa),
intenta reconèixer les parts següents d'un datagrama.
\begin{center}
\begin{bytefield}{48}
\bitbox{12}{Layer~2 header} & \bitbox{12}{IP header} & \bitbox{24}{Datagram payload}
\end{bytefield}
\end{center}
Si les reconeix, el router processa el datagrama de la manera següent:
\begin{enumerate}
\item Es descarta tot el frame $F$ tret que la MAC de destinació (a la capçalera de capa~2) 
  sigui la de la interfície del router que ha detectat el frame.\\[-0.25cm]
  
\item El router mira l'adreça IP de destinació (a la capçalera IP de $F$) i la seva \conceptRef{routing table}{taula d'encaminament} 
  per decidir el següent \conceptRef{hop}{salt} del datagrama, \ie, el dispositiu al qual aquest router passarà 
  el datagrama, i quina \conceptRef{interface}{interfície} farà servir (vegeu la secció~\ref{sec:layer3:routing} més avall).\\[-0.25cm]
  
\item El router crea i envia un frame $F'$ fent els dos canvis següents a $F$:\\[-0.25cm]
  \begin{enumerate}[label=\alph*)]
  \item La capçalera de capa~2 de $F'$ canvia completament:
    \begin{itemize} 
    \item La \concept{MAC} de destinació és la del següent salt,
    i la MAC d'origen és la del router que reenvia el datagrama.
    La IP del següent salt la dona la \conceptRef{routing table}{taula d'encaminament} (vegeu la secció~\ref{sec:layer3:routing}), 
    i la seva MAC s'obté a partir d'aquesta IP amb \concept{ARP} (secció~\ref{sec:layer3:arp}).\\[-0.4cm]
    %     
    \item La MAC d'origen del frame $F'$
    és diferent de la MAC de destinació del frame rebut $F$ si el router 
    ha rebut $F$ per una interfície (\eg \inlineCode{eth7}) i el reenvia per una altra 
    (\eg,  \inlineCode{eth9}).
    \end{itemize}
    
    \begin{center}     
    \vspace{0.15cm}
    \includegraphics[width=0.8\linewidth]{router_ports.jpg}
    \vspace{0.15cm}
    \end{center}
    
  \item La capçalera IP de $F'$ difereix de la de $F$ en el següent:
    \begin{itemize}
    \item El camp \concept{TTL} (temps de vida) es decrementa en $1$.\\
    Si el TTL arriba a $0$, el datagrama es descarta en lloc de reenviar-se.
    \item Quan el camp \conceptRef{IP Options}{Options} és present, sovint s'actualitza.
    \item El \concept{checksum} (suma de verificació) de la capçalera s'actualitza en conseqüència.\\
    \end{itemize}
  \end{enumerate}
\end{enumerate}

\begin{remark}
 Fixa't que les adreces IP i el payload es copien sense modificar
  quan es reenvia un datagrama.
\end{remark}




Quan un router rep un \conceptRef{datagram}{datagrama} IP per reenviar i aquest router 
està connectat a la mateixa LAN que l'adreça IP de destinació, el router fa un \conceptRef{direct delivery}{lliurament directe} (DD).
% 
En cas contrari, el datagrama es reenvia al següent router en un \conceptRef{indirect delivery}{lliurament indirecte} (ID).
% 
Un router sap si ha de fer DD o ID, així com l'adreça IP de destinació (i d'origen) que ha de fer servir,
gràcies a la seva \conceptRef{routing table}{taula d'encaminament}.

Totes les taules d'encaminament tenen com a mínim les columnes següents:
\begin{itemize}
\item \textbf{Bloc d'adreces de destinació}: una o més adreces IP que es poden expressar com una combinació IP + màscara de xarxa.
  Una fila de la taula només es pot aplicar a un datagrama si la seva IP de destinació és dins del bloc d'adreces de destinació d'aquella fila:
  \begin{itemize}
  \item Si només s'aplica una fila de la taula, aquesta fila és la que es fa servir per encaminar.
  \item Si s'aplica més d'una fila a un datagrama, es fa servir la fila amb el bloc d'adreces de destinació més petit 
    (el més específic).
  \end{itemize}
  Un bloc de destinació \conceptRef{default routing table}{\textit{default}} indica el bloc \otherBase{0.0.0.0/0}.
  Com que és el bloc més gran possible, només s'aplica si no s'aplica cap de les altres files. 
  L'entrada \textit{default} no és obligatòria.\\[-0.3cm]
  
\item \textbf{Gateway/via} (passarel·la): si la fila correspon a un \conceptRef{indirect delivery}{lliurament indirecte}, aquesta columna conté l'adreça IP
  del següent router. En cas contrari, aquesta columna conté \otherBase{0.0.0.0} per indicar un \conceptRef{direct delivery}{lliurament directe}.\\[-0.3cm]
  
\item \textbf{Interfície}: nom intern de la interfície de xarxa (\eg, de la targeta Ethernet) que s'ha de fer servir per lliurar el 
  datagrama. Els sistemes operatius poden fer servir la convenció que vulguin, sempre que les interfícies es puguin identificar de manera única a la taula.
  Els noms habituals comencen per \inlineCode{eth}, \inlineCode{wlan}, \inlineCode{atm}, 
  i, més recentment, \inlineCode{en} i \inlineCode{wl}.
\end{itemize}

\begin{remark}
Tots els dispositius tenen la seva pròpia taula d'encaminament, no només els routers.
\end{remark}

\begin{exercise}
A continuació es mostra la \conceptRef{routing table}{taula d'encaminament} d'un router petit.

\begin{center}
\vspace{0.2cm}
\begin{tabular}{ccc}
\toprule
\textbf{Bloc d'adreces destí} & \textbf{gateway/via} & \textbf{Interfície} \\
\toprule
\otherBase{\textit{default}} & \otherBase{192.168.0.1} & \inlineCode{eth0}  \\
\otherBase{192.168.0.0/24} & \otherBase{0.0.0.0} & \inlineCode{eth0} \\
\otherBase{172.16.0.4/32} & \otherBase{0.0.0.0} & \inlineCode{eth1} \\
\otherBase{10.49.4.0/28} & \otherBase{10.49.4.1} & \inlineCode{eth2} \\
\otherBase{10.49.4.3/32} & \otherBase{10.49.4.2} & \inlineCode{eth2} \\
\bottomrule
\end{tabular}
\vspace{0.2cm}
\end{center}

\begin{itemize}
\item Dibuixa un mapa de xarxa coherent amb la taula d'encaminament que contingui aquest router.
  Per a les interfícies del router, indica'n el nom i algunes adreces IP i de capa~2 vàlides.\\[-0.3cm]
  
\item Com reenviaria aquest router un datagrama adreçat a una adreça IP \conceptRef{public}{pública}?\\[-0.3cm]

\item Dona una adreça IP de destinació que es processaria amb la fila \otherBase{10.49.4.0/28},
i una altra que es processaria amb la fila \otherBase{10.49.4.2/32}. 
Pensa un cas d'ús en què tingui sentit tenir aquestes dues regles.
\end{itemize}
\end{exercise}

Perquè Internet funcioni, els routers que interconnecten LANs a escala global han d'estar ben configurats. 
De vegades no és així, i això pot provocar \conceptRef{routing errors}{errors d'encaminament}. Dos problemes importants 
que et pots trobar són:
\begin{itemize}
\item ``Network is unreachable'': cap de les files de la taula d'encaminament fa servir un bloc d'adreces que contingui 
  l'adreça de destinació del datagrama. 
  El router no pot saber quin ha de ser el següent salt, de manera que el datagrama no es reenvia.
  Això provoca un error ``ICMP Destination Unreachable'' (vegeu la secció~\ref{sec:layer3:icmp}).
  
\item ``Time to Live exceeded'': cada vegada que es reenvia un datagrama,
  el seu \concept{TTL} es decrementa, tal com s'ha dit a la secció~\ref{sec:layer3:ip:format}. 
  Si diversos routers estan configurats de manera que el datagrama segueix algun tipus de cicle (bucle), 
  el datagrama es reenviaria infinitament.
  Quan el \concept{TTL} arriba a 0, el datagrama es descarta i es produeix aquest error.
  Això provoca un error ``ICMP Time Exceeded'' (vegeu la secció~\ref{sec:layer3:icmp}).
\end{itemize}

\begin{exercise}
Es pot produir l'error ``Network unreachable'' si un router té una entrada \otherBase{\textit{default}}?
\end{exercise}

\begin{exercise}
Proposa un exemple de LANs interconnectades i especifica les taules d'encaminament necessàries perquè un datagrama 
segueixi un cicle i provoqui l'error ``TTL exceeded''.
\end{exercise}

\section{Divisió en subxarxes IP}\label{sec:layer3:subnetting}
\vspace{6em}
\begin{bytefield}[bitwidth=1.35em]{32}
\bitbox[]{1}{\raisebox{2.25em}{\rotatebox{90}{$\leftarrow$\otherBase{0.0.0.0}}}} &
\bitbox[]{7}{} &
\bitbox[]{1}{\raisebox{2.5em}{\rotatebox{90}{$\leftarrow$\otherBase{64.0.0.0}}}} &
\bitbox[]{7}{} &
\bitbox[]{1}{\raisebox{3.25em}{\rotatebox{90}{$\leftarrow$\otherBase{128.0.0.0}}}} & 
\bitbox[]{7}{} &
\bitbox[]{1}{\raisebox{3em}{\rotatebox{90}{$\leftarrow$\otherBase{192.0.0.0}}}}  &
\bitbox[]{6}{} &
\bitbox[]{1}{\raisebox{7em}{\rotatebox{90}{$\leftarrow$\otherBase{255.255.255.255}}}}  \\
\\
\bitbox{32}{\otherBase{0.0.0.0/0}} \\
\bitbox{16}{\otherBase{0.0.0.0/1}} & \bitbox{16}{\otherBase{128.0.0.0/1}} \\
\bitbox{8}{\otherBase{0.0.0.0/2}} & \bitbox{8}{\otherBase{64.0.0.0/2}} & \bitbox{8}{\otherBase{128.0.0.0/2}} & \bitbox{8}{\otherBase{192.0.0.0/2}} \\
% 
\bitbox{4}{\otherBase{/3}} & \bitbox{4}{\otherBase{/3}} & \bitbox{4}{\otherBase{/3}} & \bitbox{4}{\otherBase{/3}} &
\bitbox{4}{\otherBase{/3}} & \bitbox{4}{\otherBase{/3}} & \bitbox{4}{\otherBase{/3}} & \bitbox{4}{\otherBase{/3}} \\

\bitbox{2}{\otherBase{/4}} & \bitbox{2}{\otherBase{/4}} & \bitbox{2}{\otherBase{/4}} & \bitbox{2}{\otherBase{/4}} & 
\bitbox{2}{\otherBase{/4}} & \bitbox{2}{\otherBase{/4}} & \bitbox{2}{\otherBase{/4}} & \bitbox{2}{\otherBase{/4}} &
\bitbox{2}{\otherBase{/4}} & \bitbox{2}{\otherBase{/4}} & \bitbox{2}{\otherBase{/4}} & \bitbox{2}{\otherBase{/4}} & 
\bitbox{2}{\otherBase{/4}} & \bitbox{2}{\otherBase{/4}} & \bitbox{2}{\otherBase{/4}} & \bitbox{2}{\otherBase{/4}} \\ 
% 
\bitbox{1}{\raisebox{0.25em}{$\vdots$}} & \bitbox{1}{\raisebox{0.25em}{$\vdots$}} & \bitbox{1}{\raisebox{0.25em}{$\vdots$}} & \bitbox{1}{\raisebox{0.25em}{$\vdots$}} & 
\bitbox{1}{\raisebox{0.25em}{$\vdots$}} & \bitbox{1}{\raisebox{0.25em}{$\vdots$}} & \bitbox{1}{\raisebox{0.25em}{$\vdots$}} & \bitbox{1}{\raisebox{0.25em}{$\vdots$}} & 
\bitbox{1}{\raisebox{0.25em}{$\vdots$}} & \bitbox{1}{\raisebox{0.25em}{$\vdots$}} & \bitbox{1}{\raisebox{0.25em}{$\vdots$}} & \bitbox{1}{\raisebox{0.25em}{$\vdots$}} & 
\bitbox{1}{\raisebox{0.25em}{$\vdots$}} & \bitbox{1}{\raisebox{0.25em}{$\vdots$}} & \bitbox{1}{\raisebox{0.25em}{$\vdots$}} & \bitbox{1}{\raisebox{0.25em}{$\vdots$}} & 
\bitbox{1}{\raisebox{0.25em}{$\vdots$}} & \bitbox{1}{\raisebox{0.25em}{$\vdots$}} & \bitbox{1}{\raisebox{0.25em}{$\vdots$}} & \bitbox{1}{\raisebox{0.25em}{$\vdots$}} & 
\bitbox{1}{\raisebox{0.25em}{$\vdots$}} & \bitbox{1}{\raisebox{0.25em}{$\vdots$}} & \bitbox{1}{\raisebox{0.25em}{$\vdots$}} & \bitbox{1}{\raisebox{0.25em}{$\vdots$}} & 
\bitbox{1}{\raisebox{0.25em}{$\vdots$}} & \bitbox{1}{\raisebox{0.25em}{$\vdots$}} & \bitbox{1}{\raisebox{0.25em}{$\vdots$}} & \bitbox{1}{\raisebox{0.25em}{$\vdots$}} & 
\bitbox{1}{\raisebox{0.25em}{$\vdots$}} & \bitbox{1}{\raisebox{0.25em}{$\vdots$}} & \bitbox{1}{\raisebox{0.25em}{$\vdots$}} & \bitbox{1}{\raisebox{0.25em}{$\vdots$}} \\
\end{bytefield}

Les \conceptRef{LAN}{LANs} es poden considerar blocs d'adreces (\eg, expressats en format \concept{CIDR})
amb dues adreces reservades. A les xarxes accessibles directament des d'Internet se'ls ha d'assignar 
un bloc vàlid prou gran per connectar-hi tots els dispositius necessaris. Per garantir que no hi hagi dos dispositius a Internet
amb la mateixa IP \conceptRef{public}{pública}, no pot haver-hi dues LANs que continguin la mateixa adreça 
(tret que l'adreça sigui \conceptRef{private IP}{privada}).

Les $2^{32}$~adreces disponibles a \concept{IPv4} es divideixen en blocs mitjançant 
\conceptRef{netmask}{màscares de xarxa}. Els blocs més grans es poden dividir per la meitat augmentant 
en $1$ la longitud de la màscara de xarxa, com es veu a la figura anterior.
% 
Per exemple, el bloc inicial \otherBase{0.0.0.0/0} (totes les adreces) 
es pot dividir en dos blocs \otherBase{/1}, cadascun amb la meitat de les adreces \concept{IPv4} 
($2^{31}$ -- encara massa gran per a qualsevol LAN real).
% 
Aquest procés continua de manera independent per a cada bloc resultant, i s'atura quan la mida de LAN 
resultant és la desitjada. 

Donat un bloc d'adreces \otherBase{/N} (\eg, d'una única LAN existent), potser volem dividir-lo 
en subblocs i tenir, en lloc seu, diverses LANs més petites (\conceptRef{subnet}{subxarxes}). 
% 
A les subxarxes només podem fer servir adreces del bloc original, però es poden reservar blocs
per a un ús futur:

\begin{exercise}
Ets l'administrador de la xarxa \otherBase{192.168.0.0/24} d'un 
laboratori de recerca. T'agradaria dividir-la en les tres subxarxes destacades a continuació.

\begin{itemize}
\item Fes un diagrama de les tres xarxes, incloent-hi un router connectat a totes tres
  mitjançant tres interfícies diferents.\\[-0.4cm]
  
\item Dona la taula d'encaminament \textit{mínima} d'aquest router (aquestes xarxes
no estan connectades a Internet).\\[-0.4cm]
  
\item Quants dispositius es poden connectar en total a aquestes tres xarxes, incloent-hi els routers?\\[-0.4cm]

\item Quantes adreces de la xarxa original queden sense assignar?\\[-0.4cm]

\item Quina és la LAN més gran que podem crear amb les adreces sense assignar?\\[-0.75cm]
\end{itemize}
% 
\begin{center}
\begin{bytefield}[bitwidth=1.375em]{32}
\bitbox{32}{\otherBase{192.168.0.0/24}} \\
% 
\bitbox{16}{\otherBase{192.168.0.0/25}} 
& \bitbox{16}{\otherBase{192.168.0.128/25}} \\
% 
\bitbox{8}[bgcolor=color1!40!white]{\otherBase{192.168.0.0/26}} 
& \bitbox{8}{\otherBase{192.168.0.64/26}} 
& \bitbox{8}{\otherBase{192.168.0.128/26}} 
& \bitbox{8}{\otherBase{192.168.0.192/26}} \\
% 
\bitbox{4}{\otherBase{/27}} &
\bitbox{4}{\otherBase{/27}} &
\bitbox{4}{\otherBase{/27}} &
\bitbox{4}{\otherBase{/27}} &
\bitbox{4}{\otherBase{/27}} &
\bitbox{4}[bgcolor=color1!40!white]{\otherBase{/27}} &
\bitbox{4}[bgcolor=color1!40!white]{\otherBase{/27}} &
\bitbox{4}{\otherBase{/27}}
\end{bytefield}
\end{center}
\end{exercise}

El procés de crear diverses LANs més petites a partir d'un únic bloc d'adreces s'anomena \conceptRef{subnet}{divisió en subxarxes} (\textit{subnetting}).
El procés invers, combinar diverses LANs en una de més gran, s'anomena \conceptRef{supernet}{agregació en superxarxes} (\textit{supernetting}). 

Dues LANs \otherBase{/N} es poden combinar si i només si les seves adreces de xarxa 
(la seva IP sense màscara) comparteixen els primers $N-1$ bits. Aquest procés pot continuar fins arribar 
a una xarxa \otherBase{/0}, sempre que siguem propietaris de les adreces IP de les dues meitats de cada superxarxa.

\begin{exercise}\ \\[-0.5cm]
\begin{itemize}
\item Quina altra xarxa pots combinar amb \otherBase{192.168.0.0/24}
per obtenir una xarxa \otherBase{/23} $N_{/23}$? Quina és la seva adreça de xarxa?

\item $N_{/23}$ forma part d'exactament una xarxa \otherBase{/22}, $N_{22}$. 
  Enumera totes les xarxes \otherBase{/24} que hauries de tenir per combinar-les en $N_{22}$.
\end{itemize}
\end{exercise}

La \conceptRef{subnetting}{divisió en subxarxes} sovint es basa en quantes adreces IP han d'estar 
disponibles a cada subxarxa. 
Aquesta mida determina el nombre de bits de la màscara de la subxarxa, com s'ha vist a la secció~\ref{sec:layer3:lans_and_netmasks}.
Es poden afegir restriccions addicionals, \eg, fer que les adreces IP d'una subxarxa siguin les més altes,
o més baixes que les d'altres subxarxes. 

\begin{remark}
Considera per exemple la xarxa \otherBase{192.168.0.0/16}. 
Volem dividir aquesta xarxa en subxarxes i crear la xarxa més petita que pugui 
connectar $100$ dispositius simultàniament. 

Com s'ha vist a la secció~\ref{sec:layer3:netmasks}, això determina que 
si s'ha de crear una subxarxa a partir d'un  per contenir com a mínim 100 nodes,
cal dedicar com a mínim $7$~bits d'adreça ($2^6 - 2 = 62,\,2^7=126$) 
als identificadors de dispositiu i els $32-7=25$ restants 
a identificar la xarxa, \ie, la subxarxa pot ser tan petita com una \otherBase{/25}.
\ \\
\begin{center}
\begin{bytefield}[bitheight=2em]{32}
\bitheader{0-31} \\
\bitbox{16}{\otherBase{192.168}}
& \bitbox{9}{Free bits}
& \bitbox{7}{Device ID}
\end{bytefield}
\end{center}

Els $9$ bits restants (posicions $16$ a $24$) es poden fixar lliurement si no
tenim cap altra restricció. Triar un valor per a aquests bits vol dir
triar una de les $2^9$ xarxes de mida \otherBase{/25} contingudes 
a la xarxa original, \otherBase{192.168.0.0/16}.
Posar aquests bits lliures a \otherBase{0} dona el rang d'IP més baix,
i posar-los a \otherBase{1}, el més alt.

Quan s'han de crear diverses subxarxes, és essencial verificar 
que cap subxarxa no estigui continguda en una altra. Si no, la mateixa adreça IP 
pertanyeria a dues LANs, i l'encaminament seria impossible. Per exemple,
si s'han de crear tres subxarxes per a 100, 500 i 200 dispositius a partir de \otherBase{192.168.0.0/16},
amb rangs d'adreces en aquest ordre, una divisió possible seria:

\begin{center}
\begin{bytefield}{32}
\bitheader{0-31}\\
\bitbox{16}{\otherBase{192.168}} & \bitboxes{1}{0000 0000 0} & \bitbox{7}{Device ID}
& \bitbox[]{2}{$\rightarrow$} & \bitbox[]{8}{\otherBase{192.168.0.0/25}} \\
\bitbox{16}{\otherBase{192.168}} & \bitboxes{1}{0000 001} & \bitbox{9}{Device ID}
& \bitbox[]{2}{$\rightarrow$} & \bitbox[]{8}{\otherBase{192.168.2.0/23}} \\
\bitbox{16}{\otherBase{192.168}} & \bitboxes{1}{0000 0001} & \bitbox{8}{Device ID}
& \bitbox[]{2}{$\rightarrow$} & \bitbox[]{8}{\otherBase{192.168.1.0/24}} \\
\end{bytefield}
\end{center}

Recorda que l'adreça de xarxa de cada subxarxa resultant té tots els bits de l'identificador de dispositiu a \otherBase{0},
i que aquesta adreça normalment es presenta en format \concept{quad decimal}.
\end{remark}

\begin{exercise}\ \\[-0.5cm]
\begin{itemize}
\item A l'exemple anterior, per què els bits lliures de la segona subxarxa són \otherBase{0000001} 
i no \otherBase{0000000}?\\[-0.4cm]

\item Repeteix l'exercici assignant les adreces més altes possibles a la tercera subxarxa,
  després les següents adreces més altes possibles a la segona, 
  i les adreces més baixes possibles a la primera.\\[-0.4cm]
  
\item Quina és la IP de \concept{broadcast} de cada subxarxa?
\end{itemize}
\end{exercise}


% 
\section{Fragmentació IP}
\label{sec:layer3:fragmentation}

Abans que un dispositiu (incloent-hi els \conceptRef{router}{routers} i els ordinadors normals)
pugui enviar un \conceptRef{datagram}{datagrama} per una LAN, ha de verificar que la seva longitud en bytes 
no sigui més gran que la \concept{MTU} d'aquesta LAN. Si ho és, el datagrama s'ha de 
\conceptRef{fragmentation}{fragmentar} en dos o més datagrames més petits que sí que càpiguen en la MTU de la LAN.

Els fragments més petits són datagrames complets amb capçalera i payload. Per tant, 
més fragments impliquen una \conceptRef{overhead}{sobrecàrrega} més gran. Cadascun d'aquests fragments s'envia 
i es reenvia individualment com qualsevol altre datagrama, i pot ser que s'hagi de tornar a fragmentar 
si passa per una altra LAN amb una MTU encara més petita.

Quan s'ha de fragmentar un datagrama, la pila de xarxa segueix 
els passos següents:
\begin{enumerate}

\item Es comprova el flag \concept{DF} (don't fragment) del datagrama original. 
    Si DF=\otherBase{1} però cal fragmentar, el datagrama no es fragmenta 
    i es produeix un error.\\[-0.2cm]

\item Les dades del \conceptRef{payload}{payload} es divideixen en blocs de mida com a màxim $S$, de manera que $S + H = MTU$ 
  ($H$ és la longitud de la capçalera del datagrama, normalment $20$~bytes). Només l'últim bloc
  pot tenir una mida inferior a $S$. Això garanteix que cap fragment no superi la MTU,
  que és precisament l'objectiu de la fragmentació.\\[-0.2cm]
  
\item La \conceptRef{header}{capçalera} del datagrama original es copia a la dels fragments amb les modificacions següents:
  \begin{itemize}
  \item El camp Datagram ID pren el mateix valor en tots els fragments d'un datagrama. 
    El sistema operatiu s'assegura de triar un ID que la seva pila de xarxa no hagi fet servir recentment.
    El camp Datagram ID s'ignora en els datagrames no fragmentats.\\[-0.3cm]
    
  \item El camp \concept{Offset} (desplaçament) de cada fragment indica la posició (índex de byte)
    del seu primer byte de dades de payload dins del payload del datagrama abans de la fragmentació.
    Aquest camp no s'expressa en bytes, sinó en paraules de $8$~bytes.
    %     
    Donat un offset en bytes, el camp \concept{Offset} es calcula dividint per $8$.
    Per tant, els offsets dels fragments només poden ser múltiples de $8$.\\[-0.3cm]
    
  \item El flag \concept{MF} (more fragments) es posa a \otherBase{1} (activat) 
    en tots els fragments excepte l'últim. Això ha de ser cert fins i tot si un o més
    fragments es tornen a fragmentar en una LAN posterior. Així, quan es fragmenta un 
    datagrama $D$ amb MF=\otherBase{1}, cap dels fragments resultants\\[-0.2cm]
  \end{itemize}
  
\item Els datagrames fragment s'envien de la mateixa manera que s'enviaria un datagrama normal.
  Els fragments es poden tornar a fragmentar, però només els torna a acoblar el host (equip connectat a la xarxa) de destinació.
  
\end{enumerate}


\begin{exercise}
Donats l'\concept{Offset} i el flag \concept{MF}, com podem distingir 
un datagrama que no s'ha fragmentat d'un fragment de datagrama?
\end{exercise}

\begin{exercise}
Per què seria problemàtic posar MF=\otherBase{0} a l'últim datagrama 
produït en fragmentar un datagrama amb MF=\otherBase{1}?
\end{exercise}

\begin{exercise}
Un dispositiu ha d'enviar un datagrama amb una capçalera IP de $20$~bytes i un payload de $2170$~bytes.
El dispositiu està connectat a la LAN1, i el datagrama ha de travessar 
la LAN2 i la LAN3 fins que arriba a la LAN4 de destinació. Les MTU d'aquestes LANs són les següents:
\begin{center}
\begin{tabular}{r cccc}
\toprule
\textbf{LAN} & LAN1 & LAN2 & LAN3 & LAN4 \\
\textbf{MTU} & 1500 & 710  & 4096 & 864  \\
\bottomrule
\end{tabular}
\end{center}

\begin{itemize}
\item Indica l'offset i la longitud del payload (tots dos en bytes),
 així com els flags MF i DF dels $4$ datagrames
 que arriben a la LAN4 com a resultat d'enviar el datagrama
 (suposant que no hi ha pèrdues ni duplicacions).
  
\item Pot el dispositiu d'origen conèixer les MTU de les LANs intermèdies i de 
destinació?
\end{itemize}

\end{exercise}


\begin{exercise}
Els datagrames que contenen fragments poden arribar duplicats i/o desordenats. 
Com hauria de gestionar aquests datagrames la pila de xarxa receptora?
\end{exercise}


\section{ARP -- Address Resolution Protocol}\label{sec:layer3:arp}

L'Address Resolution Protocol (\concept{ARP}) el fa servir la capa~3 (\concept{IP})
per obtenir l'adreça \concept{MAC} (física) de destinació corresponent a una adreça IP coneguda.

Els paquets ARP s'envien com a payload de frames de capa~2 (\eg, Ethernet). 
Els paquets ARP \textit{no} són datagrames i \textit{no} contenen capçaleres IP.

\begin{remark}
La IP que consulta la capa~3 és la del \concept{router} per al \conceptRef{indirect delivery}{lliurament indirecte} i la del dispositiu de destinació
per al \conceptRef{indirect delivery}{lliurament indirecte}. Recorda també que la IP de destinació del datagrama no canvia durant el reenviament.
\end{remark}

\subsection{Format del paquet}

\vspace{0.2cm}
\begin{center}
\begin{bytefield}[bitheight=3em,bitwidth=1.1em]{32}
\bitheader{0-31}\\
% 
\bitbox{16}{Layer~2 address type \\ (\otherBase{1}: Ethernet)} 
& \bitbox{16}{Layer~3 address type \\ (\otherBase{0x0800}: IP)} \\ 
% 
\bitbox{8}{Layer~2\\address length} & \bitbox{8}{Layer~3\\address length} 
& \bitbox{16}{ARP operation\\(\otherBase{1}: request, \otherBase{2}: reply)} \\
% 
\bitbox[tl]{16}[bgcolor=color2!50!white]{Sender Layer~2 address} &
\bitbox[trb]{16}[bgcolor=color2!50!white]{} \\
\bitbox[lbr]{16}[bgcolor=color2!50!white]{} 
& \bitbox{16}[bgcolor=color2!50!white]{Sender Layer~3 address ...} \\
\bitbox{16}[bgcolor=color2!50!white]{... Sender Layer~3 address}
& \bitbox[ltr]{16}[bgcolor=color1!50!white]{Receiver Layer~2 address} \\
\bitbox[tlb]{16}[bgcolor=color1!50!white]{} & \bitbox[rb]{16}[bgcolor=color1!50!white]{} \\
\bitbox{32}[bgcolor=color1!50!white]{Receiver Layer~3 address}
\end{bytefield}
\end{center}

\begin{remark}
El protocol \concept{ARP} \textit{no} fa servir capçaleres IP.
\end{remark}
\begin{remark}
``Sender'' i ``receiver'' s'han d'interpretar com ``... d'aquest missatge ARP''.
\end{remark}


\begin{exercise}\ \\[-0.5cm]
\begin{itemize}
\item Quina és la finalitat dels $4$ primers camps d'un paquet ARP?
\item Quin efecte tenen en els camps restants?
\end{itemize}
\end{exercise}

\begin{exercise}
Quantes adreces vàlides apareixen en un frame Ethernet que transporta una petició ARP?
I en una resposta ARP?
\end{exercise}


\subsection{Funcionament}

Quan un node \node{A} necessita traduir una adreça IP a una adreça MAC, envia un frame de \concept{broadcast}
(MAC de destinació \otherBase{ff:ff:ff:ff:ff:ff} que conté un paquet ARP. 
En aquest paquet ARP, l'operació es posa a \inlineCode{REQUEST} (\otherBase{1}),
la IP de destinació és la que es vol traduir, i
la MAC de destinació es posa a \otherBase{00:00:00:00:00:00}.

Quan una pila de xarxa \node{B} veu un paquet de petició ARP, 
envia un paquet de resposta ARP només si 
la IP que té configurada és l'adreça IP de destinació de l'ARP.
Aquesta resposta és idèntica a la petició amb aquests canvis:
\begin{itemize}
\item El codi OP del paquet ARP es posa a \inlineCode{REPLY} (\otherBase{2}).
\item ``Origen'' ara fa referència a \node{B} (a la petició feia referència a \node{A}),
  i ``destinació'', a \node{A}.
\item En una resposta s'omplen totes $4$ adreces del paquet 
  (no poden ser \otherBase{00:00:00:00:00:00}).
\end{itemize}

\begin{exercise}
Quan s'executa l'ordre \inlineCode{ip route} en un node, la sortida conté:

\begin{verbatim}
default via 192.168.0.1 dev wlp5s0 src 192.168.0.2 metric 600 [...]
\end{verbatim}

Al seu torn, l'ordre \inlineCode{ip link} conté les línies següents:

\begin{verbatim}
2: enp42s0: <NO-CARRIER,BROADCAST,MULTICAST,UP> mtu 1500 state DOWN [...]
    link/ether d8:43:ae:61:ed:a4 brd ff:ff:ff:ff:ff:ff
3: wlp5s0: <BROADCAST,MULTICAST,UP,LOWER_UP> mtu 1500 state UP [...]
    link/ether 5c:e4:3a:18:97:1b brd ff:ff:ff:ff:ff:ff 
\end{verbatim}

\begin{itemize}
\item El node ha de fer el \conceptRef{indirect delivery}{lliurament indirecte} d'un datagrama,
però no coneix l'adreça MAC del seu router gateway. Indica el valor de 
tots els camps (incloent-hi la capçalera Ethernet i el missatge ARP) de la petició ARP 
que enviarà el node per obtenir la MAC del router.

\item Descriu els valors dels camps de la resposta ARP (pots fer servir qualsevol adreça vàlida
      com a adreça MAC del router).
\end{itemize}
\end{exercise}

No és gaire probable que l'adreça MAC associada a una IP canviï sovint.
Quan es rep una resposta ARP vàlida, el sistema operatiu actualitza una 
\conceptRef{cache}{memòria cau} afegint-hi (o actualitzant-hi) les traduccions IP $\leftrightarrow$ MAC rebudes.
% 
Si el sistema operatiu ha d'enviar un datagrama a aquesta IP abans que caduqui l'entrada de la memòria cau,
es fa servir directament l'adreça MAC de la taula sense necessitat de missatges ARP addicionals.
% 
L'ordre \inlineCode{ip neighbor} mostra el contingut d'aquesta taula, \eg,
\begin{verbatim}
192.168.0.1 dev wlp5s0 lladdr cc:2d:14:fc:22:60 REACHABLE
\end{verbatim}

\begin{remark}
ARP es pot fer servir per \conceptRef{scan}{escanejar} les IP utilitzades en una LAN. 
En particular, un node pot fer servir ARP per verificar que ningú més no fa servir la seva adreça IP configurada.
\end{remark}

ARP es va proposar en una època en què la seguretat de xarxa encara no era un problema. Com a conseqüència,
no proporciona cap mecanisme per verificar que una resposta ARP determinada sigui correcta.
% 
Això es pot aprofitar per \conceptRef{poison}{enverinar} la taula de memòria cau ARP d'un dispositiu enviant 
respostes ARP \conceptRef{gratuitous}{gratuïtes} (no sol·licitades) a la víctima amb 
associacions IP $\leftrightarrow$ MAC falses.

Els sistemes operatius moderns normalment ignoren les respostes ARP no sol·licitades, 
i els \conceptRef{switch}{switches} (commutador) moderns filtren els missatges ARP que no coincideixen amb les
adreces IP i MAC reals. Tot i així, IPv6 abandona ARP en favor del 
protocol Neighbor Discovery (ND), més complex.

\begin{exercise}
Quin impacte podria tenir enverinar la memòria cau ARP d'una víctima, 
canviant la MAC emmagatzemada del router gateway per la MAC de l'atacant?
\end{exercise}

\section{ICMP -- Internet Control Message Protocol}\label{sec:layer3:icmp}

L'Internet Control Message Protocol (\concept{ICMP}) és un protocol
auxiliar (tècnicament part d'\concept{IP}) que es fa servir per provar i depurar
les comunicacions de xarxa (i entre xarxes). 


\subsection{Format del paquet}

Els missatges ICMP s'encapsulen com a payload de datagrames IP.
A la capçalera IP, el \conceptRef{payload}{protocol del payload} es posa a \otherBase{0x01} (``ICMP'').\\[-0.2cm]

\begin{center}
\begin{bytefield}[bitwidth=1.2em]{32}
\bitheader{0-31}\\
\bitbox{32}[bitheight=3em]{IP Header (see \ref{sec:layer3:ip:format})} \\
\bitbox{8}{ICMP type} & \bitbox{8}{ICMP subtype} & \bitbox{16}{ICMP checksum} \\
\bitbox{32}{Rest of ICMP header (mandatory, depends on type)} \\
\bitbox[tlr]{32}[bitheight=3em]{ICMP data\\{\scriptsize(variable length, byte aligned)}} \\
\bitbox[lbr]{24}{\hspace{3cm}\raisebox{0.25cm}{$\vdots$}} & \bitbox[lt]{8}{} \\
\end{bytefield}
\end{center}

Les combinacions de tipus/subtipus ICMP més habituals que és probable que et trobis avui dia:
\begin{center}
\begin{tabular}{ccc l}
\toprule
\textbf{Tipus ICMP} & \textbf{Subtipus ICMP} & \textbf{Error?} & \textbf{Significat} \\
\toprule
$8$ & $0$       &              & Missatge ICMP Echo Request (\concept{ping}). \\
$0$ & $0$       &              & Missatge ICMP Echo Reply (resposta \concept{ping}). \\
$3$ & $0$--$16$ & $\checkmark$ & Destination Unreachable. Subtipus: causa. \\
$11$ & $0$      & $\checkmark$ & Time to Live Exceeded: \concept{TTL} arriba a $0$.  \\
\bottomrule
\end{tabular}
\end{center}

En aquests missatges ICMP d'error, el camp de dades ICMP conté informació sobre el 
datagrama que ha causat el missatge ICMP d'error:\\[-0.2cm]

\begin{center}
\begin{bytefield}[bitwidth=1.2em]{32}
\bitheader{0,7,8,15,16,23,24,31}\\
\bitbox{32}[bitheight=3em]{Problematic datagram's complete IP header} \\
\bitbox{32}[bitheight=2.5em]{Problematic datagram's first $8$~payload bytes} \\
\end{bytefield}
\end{center}

En els missatges echo request i echo reply, el camp de dades ICMP és arbitrari dins dels límits
dels datagrames IP. La fragmentació de datagrames ICMP (\eg, payloads de ping grans) és possible,
però molts routers ignoren els fragments ICMP.

\subsection{Funcionament}

Els missatges d'error ICMP els envien els \conceptRef{router}{routers} i els hosts normals,
desencadenats per altres datagrames:
\begin{itemize}
\item La família de missatges ``Destination unreachable'' s'envia quan un datagrama no es pot 
lliurar a la seva destinació, \eg, per una \conceptRef{routing table}{taula d'encaminament} mal configurada, 
un \concept{port} tancat (vegeu la secció~\ref{sec:layer4}), o simplement perquè el host
de destinació no està engegat.

\item El missatge ``TTL exceeded'' s'envia quan un router rep un datagrama per a 
\conceptRef{indirect delivery}{lliurament indirecte} amb \concept{TTL}$=1$ (vegeu la secció~\ref{sec:layer3:routing}).
\end{itemize}

Per evitar quantitats infinites de trànsit, \textit{no} es generen errors ICMP
quan els datagrames que els provoquen són:
\begin{itemize}
\item Datagrames ICMP (d'error o no).
\item Fragments de datagrama, excepte el primer.
\end{itemize}

\begin{exercise}
Com podem distingir el primer fragment d'un datagrama dels fragments següents?
\end{exercise}
